\chapter{The reproduction}
\label{ch:reproduction}

\section{What "reproduce" means here, and the hardware}

There are degrees of reproduction. From weakest to strongest:

\begin{enumerate}[leftmargin=*,itemsep=2pt]
  \item Re-run the authors' notebook unchanged.
  \item Re-implement the analysis from the paper, using their released model.
  \item Retrain the model from scratch and re-derive everything.
\end{enumerate}

This project is a thorough version of (2). I rebuilt the data pipeline,
the metrics and the analysis from the paper's description and primary sources,
and used the released checkpoint. I did not retrain --- pretraining on 600{,}000
proteins is well beyond a laptop.

Hardware: an NVIDIA RTX 4050 laptop GPU with 6\,GB of memory. The model is
253.6M parameters, which at half precision is about 0.5\,GB --- comfortable for
inference. Total compute for everything in these notes was roughly six
GPU-hours.

\begin{keybox}{The discipline I held to}
Every choice the paper does not specify is written down as a numbered
assumption (A1--A15) with its reasoning and, where possible, evidence. Every
place my result differs from the paper's, or where the paper is internally
inconsistent, is a numbered discrepancy (D1--D15) stated neutrally. Every script
writes a JSON manifest recording the git commit, the model revision, the random
seed, the platform and the version of every library. No number in these notes
exists without a file behind it.
\end{keybox}

\section{Step 1: is the released model the model described?}
\label{sec:checkpoint}

Before anything else: does the checkpoint on Hugging Face match the paper?

\begin{table}[htbp]
\centering\small
\caption{Architecture check. \code{scripts/00\_check\_env.py}.}
\begin{tabular}{lcc}
\toprule
& \textbf{paper} & \textbf{checkpoint} \\
\midrule
layers & 12 & 12 \\
attention heads & 8 & 8 \\
hidden size & 1024 & 1024 \\
feed-forward size & 4096 & 4096 \\
parameters & 253.6\,M & \textbf{253{,}556{,}736} \\
architecture & GPT-NeoX-style & \code{gpt\_neox} \\
\bottomrule
\end{tabular}
\end{table}

All four stated hyperparameters match and the parameter count agrees to the
paper's rounding. The resolved commit hash is recorded so a future run uses
identical weights.

\subsection{The sharpest test available}

Architecture matching is necessary but weak --- it says nothing about the
weights. Something stronger was needed: a test with a \emph{known correct
answer}.

The paper's Experimental Section prints one fine-tuning record in full: a
sequence, and the eight values paired with it. If this really is the fine-tuned
model, feeding it that sequence should return that vector.

\begin{goodbox}{It does, exactly}
\begin{center}
\code{[0.327, 0.356, 0.261, 0.287, 0.437, 0.190, 0.220, 0.301]}
\end{center}
All eight entries, matching the published vector. This exercises the tokenizer,
the prompt format, the weights and my output parser in one shot, and unlike
every other test in the project it has an externally known answer.
\end{goodbox}

There is a wrinkle here that I initially got wrong, and which turned out to be
in my favour --- it is discrepancy D8 and I tell that story in
Section~\ref{sec:d8}.

\section{Step 2: rebuilding the 1{,}033-pair dataset}
\label{sec:dataset}

The paper says its dataset was ``constructed and curated based on the silkome
dataset'' and contains 1{,}033 pairs. It does not say how. The pairs are not in
the authors' repository. So they had to be rebuilt from the Silkome release.

My first attempt joined the two Silkome files on \code{ncbi\_tax\_id}, the
species identifier present in both. It returned \textbf{zero rows}. The field
exists in both files but shares no values --- they are different numbering
schemes.

The correct key is \code{idv\_id}: the individual spider.

\begin{figure}[htbp]
\centering
\begin{tikzpicture}[font=\small,
  b/.style={draw=greyc!70,fill=boxbg,rounded corners=2pt,inner sep=6pt,
            text width=4.6cm,align=center},
  ar/.style={-{Stealth[length=2.4mm]},thick,greyc!80},
  num/.style={font=\scriptsize\bfseries,text=ourscol},
]
\node[b] (fasta) {\textbf{protein FASTA}\\ \scriptsize 11{,}155 records};
\node[b,right=2.6cm of fasta] (mech) {\textbf{mechanical properties}\\ \scriptsize 446 individuals};
\node[b,below=1.0cm of $(fasta)!0.5!(mech)$] (join) {\textbf{join on \code{idv\_id}}\\ \scriptsize 3{,}778 rows};
\node[b,below=0.8cm of join] (comp) {\textbf{all 8 properties present}\\ \scriptsize 3{,}563 rows};
\node[b,below=0.8cm of comp,draw=ourscol!70,fill=ourscol!8] (masp)
  {\textbf{type begins ``MaSp''}\\ \scriptsize \textbf{1{,}033 rows}};

\draw[ar] (fasta.south) -- (join.north west);
\draw[ar] (mech.south) -- (join.north east);
\draw[ar] (join) -- (comp);
\draw[ar] (comp) -- (masp);

\node[right=0.35cm of masp,text width=3.6cm,font=\scriptsize\itshape,text=ourscol]
  {the paper states 1{,}033. Exact match.};
\node[left=0.35cm of join,text width=2.9cm,font=\scriptsize\itshape,text=greyc]
  {joining on \code{ncbi\_tax\_id} instead returns \textbf{0}};
\end{tikzpicture}
\caption[Dataset reconstruction]{Rebuilding the fine-tuning set. The rule was
found by searching for one that lands on the published count; hitting 1{,}033
exactly is strong evidence it is the authors' rule, though it remains a
reconstruction (assumption A1).}
\label{fig:dataset}
\end{figure}

Three properties of the result are worth knowing:

\begin{itemize}[leftmargin=*,itemsep=3pt]
  \item \textbf{1{,}033 rows but 1{,}028 distinct sequences.} Five sequences
        appear twice with \emph{different} property values, because the same
        protein was recovered from two individuals whose fibres tested
        differently. So the dataset contains contradictory labels for five
        inputs. I kept both, since the published count requires it (A12).
  \item \textbf{553 CTD, 480 NTD.} Terminal-domain-anchored records, as
        discussed in Section~\ref{sec:silkome}.
  \item Lengths run 115--1854, median 378. Types: MaSp1 349, MaSp2 331,
        MaSp 223, MaSp3B 47, MaSp3 43, MaSp2B 40.
\end{itemize}

\section{Step 3: recovering the normalisation constants}
\label{sec:normalisation}

To convert raw measurements into the model's $[0,1]$ scale you need the eight
$(\min, \max)$ pairs. They live in the paper's Table S5, in Supporting
Information that --- at this stage of the project --- I could not obtain.

So I derived them. The Experimental Section gives one clue: the extremes were
taken ``across the entire dataset (not limited to MaSp sequences)''. That
sentence has at least three readings, so I tested all three against the one
published normalised vector I had, from the worked example.

\begin{table}[htbp]
\centering\small
\caption[Normalisation candidates]{Testing three readings of ``the entire
dataset'' against the paper's own published vector.}
\begin{tabular}{lc}
\toprule
\textbf{population used for min/max} & \textbf{max abs.\ error vs published} \\
\midrule
raw \code{mechanical\_properties.csv} (446 rows) & 0.139 \\
\textbf{sequence-joined, all spidroin types (3{,}563 rows)} & \textbf{0.00047} \\
sequence-joined, MaSp only (1{,}033 rows) & 0.054 \\
\bottomrule
\end{tabular}
\end{table}

The middle reading reproduces all eight published values to within the paper's
own three-decimal rounding.

\begin{goodbox}{Later confirmed outright}
Much later I located the Supporting Information appended to the arXiv preprint
of the same paper. Table~S5 is there, with units. Every one of my recovered
constants matches:

\begin{center}
\small
\begin{tabular}{lccc}
\toprule
property & unit & Table~S5 & recovered \\
\midrule
toughness & GJ/m$^3$ & 0.005 -- 0.39 & 0.005 -- 0.39 \\
SD toughness & GJ/m$^3$ & 0.001 -- 0.136 & 0.001 -- 0.136 \\
$E$ & GPa & 0.38 -- 37.0 & 0.38 -- 37.0 \\
SD $E$ & GPa & 0.03 -- 9.76 & 0.03 -- 9.76 \\
strength & GPa & 0.17 -- 3.33 & 0.17 -- 3.33 \\
SD strength & GPa & 0.01 -- 0.8 & 0.01 -- 0.8 \\
strain & \% & 5.1 -- 53.2 & 5.1 -- 53.2 \\
SD strain & \% & 0.1 -- 13.7 & 0.1 -- 13.7 \\
\bottomrule
\end{tabular}
\end{center}

This also confirms the dataset reconstruction from a second direction: the
constants are a property of the population they were computed over, so getting
them exactly right means I assembled the same population the authors did.
\end{goodbox}

\section{Step 4: a decoding decision, and why it needed measuring}
\label{sec:decodingchoice}

The released notebook decodes the forward task by \emph{sampling} at
temperature 0.01. I decode \emph{greedily}. That is a deliberate departure, and
since departures need justifying, I measured the consequences on 60 real silk
sequences.

\begin{table}[htbp]
\centering\small
\caption[Decoding diagnostics]{\code{scripts/01b\_forward\_decoding\_diagnostics.py}}
\begin{tabular}{lc}
\toprule
\textbf{question} & \textbf{answer} \\
\midrule
Do greedy and sampled agree? & identical on \textbf{58/60} (96.7\%) \\
Where they differ, by how much? & at most 0.38 in one slot, 2 sequences \\
Does greedy parse at the 64-token budget? & \textbf{60/60} \\
Does 32 tokens suffice? & \textbf{0/60} --- the answer is cut off \\
Is greedy invariant to batch size? & \textbf{16/16} identical \\
Is sampled invariant to batch size? & \textbf{no} \\
\bottomrule
\end{tabular}
\end{table}

The last row is the reason. I batch the forward pass for speed, and all rows of
a batch draw from one random number stream. Under sampling, a candidate's
predicted properties would depend on \emph{which other candidates happened to
share its batch} --- a reproducibility disaster that is entirely mine to solve,
since the authors did not batch.

There was also an incidental discovery: under sampling, the model sometimes
continues the amino-acid sequence \emph{before} emitting the property vector,
and can then be cut off mid-answer by the 64-token budget. The released
notebook catches that in a bare \code{except:}, so those cases silently vanish.
Under greedy decoding it never happens.

\begin{warnbox}{This choice was later exonerated}
A reader could reasonably worry that greedy decoding is why my numbers came out
lower than the paper's. I tested exactly that in
Section~\ref{sec:stochasticity} and the effect is $+0.0065$ --- two orders of
magnitude too small. Recorded as assumption A5.
\end{warnbox}

\section{Step 5: re-running the whole procedure}
\label{sec:table1}

With everything rebuilt, I ran the cycle-consistency procedure of
Figure~\ref{fig:cycle} for all eight targets, with $N = 2048$ generation
attempts each, matching the notebook's budget. About 4.5 GPU-hours.

\begin{table}[htbp]
\centering\small
\caption[Table 1 reproduction]{The headline reproduction. ``Pool'' is the
number of candidates that parsed \emph{and} were novel, and so could be scored.}
\label{tab:repro}
\begin{tabular}{lccccc}
\toprule
\textbf{set} & \textbf{paper \Rsq} & \textbf{our max} & \textbf{our mean} & \textbf{our median} & \textbf{pool} \\
\midrule
F1 & 0.8764 & 0.7221 & $-0.36$ & $-0.03$ & 815 \\
F2 & 0.8369 & 0.6964 & $-3.57$ & $-1.41$ & 739 \\
F3 & 0.6889 & \textbf{0.6780} & $-0.36$ & $+0.09$ & 735 \\
S1 & 0.8899 & 0.7120 & $-6.54$ & $-2.25$ & 712 \\
S2 & 0.5640 & 0.2234 & $-0.71$ & $-0.67$ & 738 \\
S3 & 0.7167 & 0.6494 & $-0.37$ & $-0.31$ & 752 \\
S4 & 0.7843 & 0.3487 & $-0.49$ & $-0.48$ & 710 \\
S5 & 0.7751 & 0.6233 & $-11.59$ & $-3.61$ & 641 \\
\bottomrule
\end{tabular}
\end{table}

\begin{warnbox}{Partial reproduction}
I do not reach the published value on \emph{any} of the eight sets. The gap
runs from 0.011 (F3, essentially reproduced) to 0.436 (S4), mean 0.27.
Across all eight, none of roughly 700--800 candidates matched or exceeded the
published figure.

The mechanism works --- the pipeline does produce sequences whose predicted
properties approach the request --- but the magnitude I obtain is lower. That is
a partial reproduction and I describe it as one.
\end{warnbox}

Two things in this table deserve attention beyond the gap.

\paragraph{The mean and median are terrible.} For six of eight sets, the
\emph{typical} candidate has negative \Rsq{} --- worse than guessing the average.
S5's mean is $-11.59$. The maximum is doing enormous work.

\paragraph{The pool is much smaller than the budget.} 2{,}048 attempts yield
only $\sim$700--800 scoreable candidates. Where did the rest go? About 5\% fail
to parse. The rest --- the large majority of the loss --- were discarded for not
being novel. That deserves its own measurement, in
Section~\ref{sec:regurgitation}.

\subsection{Testing one explanation for the gap}
\label{sec:stochasticity}

Before attributing the shortfall to anything, one purely statistical mechanism
had to be ruled out.

The notebook scores candidates with a \emph{sampled} forward pass. So each
candidate's \Rsq{} is itself a random variable, and a maximum over $N$ noisy
scores is inflated relative to a maximum over $N$ noiseless ones --- roughly by
the scale of the scoring noise. My scores are noiseless. Could that alone
explain 0.27?

I re-scored the \emph{identical} candidate sequences with the notebook's
sampled decoder --- same sequences, same targets, only the scorer changed.

\begin{table}[htbp]
\centering\small
\caption[Scoring stochasticity]{\code{scripts/10b\_scoring\_stochasticity.py}}
\begin{tabular}{lcccc}
\toprule
\textbf{set} & \textbf{max, greedy} & \textbf{max, sampled} & \textbf{inflation} & \textbf{gap remaining} \\
\midrule
S1 & 0.7120 & 0.7379 & $+0.026$ & 0.152 \\
F1 & 0.7221 & 0.7221 & $0.000$ & 0.154 \\
S4 & 0.3487 & 0.3487 & $0.000$ & 0.436 \\
S2 & 0.2234 & 0.2234 & $0.000$ & 0.341 \\
\bottomrule
\end{tabular}
\end{table}

Mean inflation $+0.0065$; scores identical on 97--99\% of candidates. The
hypothesis is \textbf{rejected}. At this point in the project I recorded the gap
as unexplained rather than attributing it to anything --- which was the right
call, because the answer arrived later and from an unexpected direction.

\section{Step 6: the answer, once the Supporting Information turned up}
\label{sec:exactmatch}

Late in the project I found the supplementary tables appended to the arXiv
preprint of the same paper (\code{arXiv:2309.10170}, pages 31--38). Wiley's copy
refuses automated requests, but arXiv does not.

Table~S1 lists the five sequences the paper actually generated. Table~S3 lists,
for each, the property vector its forward task predicted and the resulting
\Rsq. So both the input and the expected output are known, and the scorer can
be tested with no sampling and no search at all.

Extracting those sequences correctly took two attempts and is the subject of
Section~\ref{sec:tables1}. With them verified:

\begin{goodbox}{The forward task reproduces exactly}
\begin{center}
\small
\begin{tabular}{lccc}
\toprule
\textbf{set} & \textbf{paper \Rsq} & \textbf{ours} & \textbf{predicted vector vs Table S3} \\
\midrule
S1 & 0.8899 & \textbf{0.8899} & \textbf{8/8 slots identical} \\
S2 & 0.5640 & 0.5176 & 6/8 \\
S3 & 0.7167 & \textbf{0.7167} & \textbf{8/8} \\
S4 & 0.7843 & \textbf{0.7843} & \textbf{8/8} \\
S5 & 0.7751 & \textbf{0.7751} & \textbf{8/8} \\
\bottomrule
\end{tabular}
\end{center}

Four of five reproduce to four decimal places --- and not merely in \Rsq, but in
the entire eight-value predicted vector, slot for slot. S2 agrees on its first
six slots and differs on the last two, consistent with the paper's sampled
decoding diverging from greedy on a low-confidence tail.
\end{goodbox}

\begin{keybox}{What this settles}
The checkpoint, the prompt format, the parser, the normalisation and my
decoding choice are all \emph{exonerated}. Given the paper's own sequences, the
scorer returns the paper's own answers.

So the shortfall in Section~\ref{sec:table1} is located precisely: it is in the
\textbf{inverse task's search}. My 2{,}048 candidates simply did not contain
sequences as good as the ones the authors found. Their generated sequences are
also markedly longer --- 1096, 926, 312, 869 and 521 residues against my 530,
493, 471, 185 and 256 --- so the two searches explored different regions of the
space. Why theirs found longer and better candidates, I cannot say from the
information available.

This is a considerably stronger reproduction result than
Section~\ref{sec:table1} suggests on its own, and it is why these notes present
the two together.
\end{keybox}

\subsection{Figure 7, against the paper's own specimens}

With Table~S1 in hand, the motif analysis could finally be run the way the
paper ran it: each generated sequence against \emph{its own} two BLAST
relatives, using the Table~S4 motif definitions.

\begin{center}
\small
\begin{tabular}{lcl}
\toprule
\textbf{set} & \textbf{motif-density correlation} & \textbf{the paper's wording} \\
\midrule
S1 & 0.658 & ``particularly evident'' \\
S2 & \textbf{0.448} & ``congruence is diminished in set 2'' \\
S3 & \textbf{0.941} & ``particularly evident'' \\
S4 & \textbf{0.944} & ``particularly evident'' \\
S5 & 0.660 & --- \\
\bottomrule
\end{tabular}
\end{center}

This supports the paper's qualitative claim: set 2 is the weakest by a clear
margin, exactly as reported, and sets 3 and 4 are strongest. Set 1, which the
paper groups with 3 and 4, comes out middling here.

\begin{warnbox}{A lesson about substitutes}
Before I had Table~S1, I ran the same analysis using \emph{nearest database
neighbours} as stand-ins for the BLAST relatives. That version gave S2 $= 0.99$
and S4 $= -0.02$ --- close to the reverse of the correct ordering.

A reproduction that had stopped at the substitute specimens would have reported
something wrong, with no internal sign of trouble. Substituting one input for
another is sometimes unavoidable, but the substitution is a claim in its own
right and deserves the same scrutiny as the result.
\end{warnbox}

\subsection{Descriptors, checked against the paper's own numbers}

Table~S3 also publishes molecular weight, instability index and isoelectric
point for all fifteen sequences. Recomputing them is an independent check of
the Figure 5 pipeline:

\begin{center}
\small
\begin{tabular}{lc}
\toprule
molecular weight & \textbf{15/15} match \\
isoelectric point & \textbf{15/15} match \\
instability index & 13/15 match \\
\bottomrule
\end{tabular}
\end{center}

The two instability-index misses are both explained. Two sequences contain
\aaseq{X} --- an ``unknown residue'' code that appears in real database records.
Biopython refuses to weigh a sequence containing \aaseq{X}, yet the paper
reports a weight. Recovering the implied mass,
\[
  \frac{\text{published MW} - \text{MW with \aaseq{X} dropped}}{\text{number of \aaseq{X}}}
  = 111.985 \text{ Da}
\]
for both sequences independently --- the conventional average residue mass. So
the paper's tool weighed unknowns as average residues. Applying that convention
reproduces both weights exactly. The instability index still differs slightly
(28.88 vs 28.55) because that index is computed from \emph{dipeptides}, and
removing an \aaseq{X} joins its two neighbours into a pair that was never in the
sequence. A consequence of the unknown-residue handling, not a disagreement
about the sequence --- reported rather than tuned away.
